% exam1-prop-logic.tex (slides)

\begin{frame}[allowframebreaks]{测验 1：命题逻辑}

\begin{problem}[命题逻辑的自然推理系统]
  请使用自然推理规则证明以下结论:

  \medskip
  (4)

  \[
    \vdash ((p \to q) \to q) \to ((q \to p) \to p)
  \]
\end{problem}

\begin{logicproof}{4}
  \lnot p & assumption \\
  \begin{subproof}
    p & assumption \\
    \bot & $\lnot\elim$ 1, 2 \\
    q & $\bot\elim$ 3
  \end{subproof}
  p \to q & $\to\intro$ 2--4 \\
  (p \to q) \to q & assumption \\
  q & $\to\elim$ 5, 6 \\
  p & $\to\elim$ 1, 7 \\
  \bot & $\lnot\elim$ 1, 8 \\
  \lnot\lnot p & $\lnot\intro$ 1--9 \\
  p & $\lnot\lnot\elim$ 10 \\
  (q \to p) \to p & $\to\intro$ 2--11 \\
  ((p \to q) \to q) \to ((q \to p) \to p) & $\to\intro$ 6--12
\end{logicproof}

\medskip
(3) (德摩根律)

\[\lnot(\alpha \land \beta) \vdash (\lnot\alpha \lor \lnot\beta)\]

提示: 你可能需要多次用到排中律 $\frac{}{\alpha \lor \lnot\alpha}$。

\begin{logicproof}{4}
  \lnot(\alpha \land \beta) & premise \\
  \alpha \lor \lnot\alpha & LEM \\
  \begin{subproof}
    \alpha & assumption \\
    \beta \lor \lnot\beta & LEM \\
    \begin{subproof}
      \beta & assumption \\
      \alpha \land \beta & $\land\intro$ 3, 5 \\
      \bot & $\lnot\elim$ 1, 6 \\
      \lnot\alpha \lor \lnot\beta & $\bot\elim$ 7
    \end{subproof}
    \begin{subproof}
      \lnot\beta & assumption \\
      \lnot\alpha \lor \lnot\beta & $\lor\introright$ 9
    \end{subproof}
    \lnot\alpha \lor \lnot\beta & $\lor\elim$ 4, 5--8, 9--10
  \end{subproof}
  \begin{subproof}
    \lnot\alpha & assumption \\
    \lnot\alpha \lor \lnot\beta & $\lor\introleft$ 12
  \end{subproof}
  \lnot\alpha \lor \lnot\beta & $\lor\elim$ 2, 3--11, 12--13
\end{logicproof}

\begin{problem}[另一个命题逻辑推理系统]
  考虑一个``希尔伯特''式的命题逻辑推理系统,
  它包含三个公理 (A1)、(A2) 与 (A3):
  \begin{enumerate}
    \item[(A1)]
      \[\alpha \to (\beta \to \alpha)\]
    \item[(A2)]
      \[ (\alpha \to (\beta \to \gamma)) \to ((\alpha \to \beta) \to (\alpha \to \gamma))\]
    \item[(A3)]
      \[ (\lnot \beta \to \lnot \alpha) \to ((\lnot \beta \to \alpha) \to \beta)\]
  \end{enumerate}
  除此之外, 它还包含一个推理规则 MP (分离规则).

  在该推理系统中, 公式 $\alpha \to \alpha$ 的一种证明方式请参见讲义。
\end{problem}

\begin{solution}
  \begin{enumerate}
    \item A1 (取 $\beta=(\alpha\to\alpha)$): $\alpha\to((\alpha\to\alpha)\to\alpha)$.
    \item A2 (取 $\beta=(\alpha\to\alpha),\,\gamma=\alpha$): $(\alpha\to((\alpha\to\alpha)\to\alpha))\to((\alpha\to(\alpha\to\alpha))\to(\alpha\to\alpha))$.
    \item 由 1、2 及 MP 得 $(\alpha\to(\alpha\to\alpha))\to(\alpha\to\alpha)$。
    \item A1 (取 $\beta=\alpha$): $\alpha\to(\alpha\to\alpha)$。
    \item 由 3、4 及 MP 得 $\alpha\to\alpha$。
  \end{enumerate}
\end{solution}

\end{frame}
